% Body prose only --- no \section command in this file.
%
% Always present. Four moves, in this order:
%   1. what this document is, and which delivery it accompanies
%   2. its purpose --- a short bulleted list
%   3. what it covers, cross-referenced to the sections
%   4. what it explicitly does NOT cover
%   5. a closing statement that the delivery is a fixed artefact
%
% Shape inherited from the FAO release notes. The "does not cover" list
% is the one that earns its keep: a reader will otherwise assume that
% anything plausibly related is in scope.

This Delivery Note accompanies \texttt{\deliveryid}, [one sentence on what the delivery is].

The purpose of this note is to:
\begin{itemize}
    \item [...];
    \item [...]; and
    \item set out what is \emph{not} in the delivery, and what is not yet known about it.
\end{itemize}

\noindent This note covers:
\begin{itemize}
    \item [... (Section~\ref{sec:files})].
\end{itemize}

\noindent This note explicitly does \emph{not} cover:
\begin{itemize}
    \item \textbf{any evaluation.} [State this plainly if nothing has been scored.]
    \item [...].
\end{itemize}

\noindent The delivery is a fixed artefact. It is identified by the identifier and checksum on the
title page --- quote both --- and is not revised in place: any correction is issued as a new
delivery with a new identifier and a new checksum.
